% GENERATED FILE. Edit canonical lessons, then run npm run generate:books.
% canonical-source-hash: fnv1a64:d2e0e75ddd290bf8
% canonical-lessons: ZH-C276-mianshi, ZH-C276-cizhi, ZH-C276-jihui, ZH-C276-tuixiu, ZH-C276-dashi

\chapter{A New Job}
\label{ch:zh-a-new-job}
\hlchaptermodality{276}

\begin{chapteropening}
\textbf{By the end of this chapter:} \emph{I can say I went for an interview or quit a job, talk about a chance and retiring, and say something is a big deal.}

Complete the last lesson of chapter 276: I can say I went for an interview or quit a job, talk about a chance and retiring, and say something is a big deal.
\end{chapteropening}

\section[\texorpdfstring{\zh{面试} (miànshì)}{面试 (miànshì)}]{\zh{面试} (miànshì) — to have an interview (for a job); interview}

\label{lesson:ZH-C276-mianshi}



\begin{quote}
Before the new one: say the Mandarin for to accept, then the Mandarin for as you like.
\end{quote}

\subsection*{You'll want to know: \zh{面试}}
\textbf{\zh{面试}} — \emph{miànshì} — \textquotedblleft{}to have an interview (for a job); interview\textquotedblright{}.

\begin{quote}
\textbf{\zh{上个月我老公去面试了}} \emph{shàng ge yuè wǒ lǎogōng qù miànshì le} — My husband went for a job interview last month.
\end{quote}

\begin{cousinweb}
\textbf{\zh{面试}} (\emph{miànshì}) is built from \zh{面} 'the face' and \zh{试} 'a test', both inherited from Middle Chinese, \emph{mjienH} and \emph{syiH} (Guangyun, in Baxter's transcription), which glosses the first 'to face' and the second 'to use'.
\end{cousinweb}

\begin{grammarlens}[title={What you've built}]
One more word for talking about the past.
\end{grammarlens}

\subsection*{Your turn}
\begin{itemize}
\raggedright
  \item \emph{Say it:} \emph{miànshì}
  \item \emph{Say it:} \emph{miànshì}, once more
  \item \emph{Say it:} \emph{shàng ge yuè wǒ lǎogōng qù miànshì le} — My husband went for a job interview last month.
  \item \emph{Recall:} say the Mandarin for to accept, then the Mandarin for as you like, then say \emph{miànshì} again
\end{itemize}

\begin{tcolorbox}[breakable,colback=teal!4,colframe=teal!35!black,title={Before you move on}]
What does \zh{面试} mean? (\textquotedblleft{}To have an interview (for a job); interview\textquotedblright{}.) Say it once more.
\end{tcolorbox}
\section[cízhí]{cízhí — to resign; to quit a job}

\label{lesson:ZH-C276-cizhi}



\begin{quote}
Before the new one: say the Mandarin for as you like, then the Mandarin for to have a job interview.
\end{quote}

\subsection*{You'll want to know: cízhí}
\emph{cízhí} — \textquotedblleft{}to resign; to quit a job\textquotedblright{}.

\begin{quote}
\emph{tā shàng ge yuè cízhí le} — She quit her job last month.
\end{quote}

Its characters come later; for now it is a word for the ear.

\begin{cousinweb}
\emph{cízhí} is built from \emph{cí} and \emph{zhí} 'a post', both inherited from Middle Chinese, \emph{zi} and \emph{tsyik} (Guangyun, in Baxter's transcription), which glosses the first 'words' and 'a plea'. The sense 'to take leave' is not given there, so how it came is not settled.
\end{cousinweb}

\begin{grammarlens}[title={What you've built}]
One more word for talking about the past.
\end{grammarlens}

\subsection*{Your turn}
\begin{itemize}
\raggedright
  \item \emph{Say it:} \emph{cízhí}
  \item \emph{Say it:} \emph{cízhí}, once more
  \item \emph{Say it:} \emph{tā shàng ge yuè cízhí le} — She quit her job last month.
  \item \emph{Recall:} say the Mandarin for as you like, then the Mandarin for to have a job interview, then say \emph{cízhí} again
\end{itemize}

\begin{tcolorbox}[breakable,colback=teal!4,colframe=teal!35!black,title={Before you move on}]
What does cízhí mean? (\textquotedblleft{}To resign; to quit a job\textquotedblright{}.) Say it once more.
\end{tcolorbox}
\section[\texorpdfstring{\zh{机会} (jīhuì)}{机会 (jīhuì)}]{\zh{机会} (jīhuì) — chance; opportunity}

\label{lesson:ZH-C276-jihui}



\begin{quote}
Before the new one: say the Mandarin for to have a job interview, then the Mandarin for to quit a job.
\end{quote}

\subsection*{You'll want to know: \zh{机会}}
\textbf{\zh{机会}} — \emph{jīhuì} — \textquotedblleft{}chance; opportunity\textquotedblright{}.

\begin{quote}
\textbf{\zh{这是一个好机会}} \emph{zhè shì yí ge hǎo jīhuì} — This is a good opportunity.
\end{quote}

\begin{cousinweb}
\textbf{\zh{机会}} (\emph{jīhuì}) is built from \zh{机} and \zh{会} 'to meet', both inherited from Middle Chinese, \emph{kj+j} and \emph{hwajH} (Guangyun, in Baxter's transcription), which glosses the first 'a meeting' and 'a trigger', and the second 'to come together'.
\end{cousinweb}

\begin{grammarlens}[title={What you've built}]
One more word for talking about the past.
\end{grammarlens}

\subsection*{Your turn}
\begin{itemize}
\raggedright
  \item \emph{Say it:} \emph{jīhuì}
  \item \emph{Say it:} \emph{jīhuì}, once more
  \item \emph{Say it:} \emph{zhè shì yí ge hǎo jīhuì} — This is a good opportunity.
  \item \emph{Recall:} say the Mandarin for to have a job interview, then the Mandarin for to quit a job, then say \emph{jīhuì} again
\end{itemize}

\begin{tcolorbox}[breakable,colback=teal!4,colframe=teal!35!black,title={Before you move on}]
What does \zh{机会} mean? (\textquotedblleft{}Chance; opportunity\textquotedblright{}.) Say it once more.
\end{tcolorbox}
\section[tuìxiū]{tuìxiū — to retire}

\label{lesson:ZH-C276-tuixiu}



\begin{quote}
Before the new one: say the Mandarin for to quit a job, then the Mandarin for an opportunity.
\end{quote}

\subsection*{You'll want to know: tuìxiū}
\emph{tuìxiū} — \textquotedblleft{}to retire\textquotedblright{}.

\begin{quote}
\emph{wǒ yéye tuìxiū le} — My grandpa has retired.
\end{quote}

Its characters come later; for now it is a word for the ear.

\begin{cousinweb}
\emph{tuìxiū} is built from \emph{tuì} 'to step back' and \emph{xiū} 'to rest', both inherited from Middle Chinese, \emph{thwojH} and \emph{xjuw} (Guangyun, in Baxter's transcription), which glosses the first 'to step back' and the second 'to rest'.
\end{cousinweb}

\begin{grammarlens}[title={What you've built}]
One more word for talking about the past.
\end{grammarlens}

\subsection*{Your turn}
\begin{itemize}
\raggedright
  \item \emph{Say it:} \emph{tuìxiū}
  \item \emph{Say it:} \emph{tuìxiū}, once more
  \item \emph{Say it:} \emph{wǒ yéye tuìxiū le} — My grandpa has retired.
  \item \emph{Recall:} say the Mandarin for to quit a job, then the Mandarin for an opportunity, then say \emph{tuìxiū} again
\end{itemize}

\begin{tcolorbox}[breakable,colback=teal!4,colframe=teal!35!black,title={Before you move on}]
What does tuìxiū mean? (\textquotedblleft{}To retire\textquotedblright{}.) Say it once more.
\end{tcolorbox}
\section[\texorpdfstring{\zh{大事} (dàshì)}{大事 (dàshì)}]{\zh{大事} (dàshì) — big deal; important matter}

\label{lesson:ZH-C276-dashi}



\begin{quote}
Before the new one: say the Mandarin for an opportunity, then the Mandarin for to retire.
\end{quote}

\subsection*{You'll want to know: \zh{大事}}
\textbf{\zh{大事}} — \emph{dàshì} — \textquotedblleft{}big deal; important matter\textquotedblright{}.

\begin{quote}
\textbf{\zh{找工作是大事}} \emph{zhǎo gōngzuò shì dàshì} — Finding a job is a big deal.
\end{quote}

\begin{cousinweb}
\textbf{\zh{大事}} (\emph{dàshì}) is built from \zh{大} 'big' and \zh{事} 'an affair', both inherited from Middle Chinese, \emph{daH} and \emph{dzriH} (Guangyun, in Baxter's transcription).
\end{cousinweb}

\begin{grammarlens}[title={What you've built}]
One more word for talking about the past.
\end{grammarlens}

\subsection*{Your turn}
\begin{itemize}
\raggedright
  \item \emph{Say it:} \emph{dàshì}
  \item \emph{Say it:} \emph{dàshì}, once more
  \item \emph{Say it:} \emph{zhǎo gōngzuò shì dàshì} — Finding a job is a big deal.
  \item \emph{Recall:} say the Mandarin for an opportunity, then the Mandarin for to retire, then say \emph{dàshì} again
\end{itemize}

\begin{tcolorbox}[breakable,colback=teal!4,colframe=teal!35!black,title={Before you move on}]
What does \zh{大事} mean? (\textquotedblleft{}Big deal; important matter\textquotedblright{}.) Say it once more.
\end{tcolorbox}
